\chapter{Axial Symmetry}\label{ch:g6:symmetry}

Fold a sheet of paper along a line: every point lands on another point,
its mirror image. This folding is \emph{axial symmetry} (or
reflection), the first geometric transformation of the book --- the
second, central symmetry, comes in \cref{ch:g7:central}.

\section{Reflections}

\begin{definition}[Symmetric point]\label{def:g6:symmetry:def}
Let $d$ be a line. The \emph{symmetric}\index{axial symmetry} of a
point $M$ across $d$ is the point $M'$ such that $d$ is the
perpendicular bisector of $[MM']$: the segment $[MM']$ is perpendicular
to $d$, and $d$ cuts it at its midpoint. A point \emph{on} $d$ is its
own symmetric.
\end{definition}

\begin{omfigure}
\begin{tikzpicture}[scale=1.05]
  \draw[thick, omExo] (0,-1.7) -- (0,1.9) node[above, font=\small] {$d$};
  \fill (-2.1,1.0) circle (1.6pt) node[left, font=\small] {$M$};
  \fill (2.1,1.0) circle (1.6pt) node[right, font=\small] {$M'$};
  \draw[dashed] (-2.1,1.0) -- (2.1,1.0);
  \draw (-0.10,0.90) -- (0.10,1.10);
  \draw (-1.15,0.90) -- (-0.95,1.10);
  \draw (0.95,0.90) -- (1.15,1.10);
  \draw[thin] (0,1.0) ++(0.16,0) -- ++(0,0.16) -- ++(-0.16,0);
  \fill (-1.3,-1.1) circle (1.6pt) node[left, font=\small] {$N$};
  \fill (1.3,-1.1) circle (1.6pt) node[right, font=\small] {$N'$};
  \draw[dashed] (-1.3,-1.1) -- (1.3,-1.1);
  \fill (0,0.2) circle (1.6pt) node[above right, font=\small] {$P = P'$};
\end{tikzpicture}

{\small Reflecting across the line $d$: the axis is perpendicular to
$[MM']$ and passes through its midpoint (equal tick marks). A point $P$
on the axis does not move.}
\end{omfigure}

\begin{method}[Constructing a symmetric point]\label{met:g6:symmetry:construct}
To construct the symmetric $M'$ of $M$ across the line $d$:
\begin{enumerate}
  \item draw the perpendicular to $d$ through $M$ (set square); call
        $H$ the point where it crosses $d$;
  \item measure $MH$ with the compass or ruler;
  \item place $M'$ on the same perpendicular, on the other side of $d$,
        with $HM' = MH$.
\end{enumerate}
On grid paper it is faster: count the squares from $M$ to the axis, and
count the same number on the other side.
\end{method}

\begin{example}[On a grid]\label{ex:g6:symmetry:grid}
If the axis is a vertical grid line and $M$ is $3$ squares to its left,
then $M'$ is $3$ squares to its right, on the same row. For a diagonal
axis, the counting works along diagonals --- always perpendicular to
the axis.
\end{example}

\begin{omfigure}
\begin{tikzpicture}[scale=0.55]
  \draw[gray!40, thin] (0,0) grid (12,6);
  \draw[thick, omExo] (6,-0.4) -- (6,6.4);
  \draw[thick, omDef, fill=omDef!15]
    (1,1) -- (4,1) -- (4,3) -- (3,3) -- (3,5) -- (1,5) -- cycle;
  \draw[thick, omThm, fill=omThm!15]
    (11,1) -- (8,1) -- (8,3) -- (9,3) -- (9,5) -- (11,5) -- cycle;
  \node[omDef, font=\small] at (2.4,0.45) {figure};
  \node[omThm, font=\small] at (9.6,0.45) {its mirror image};
\end{tikzpicture}

{\small A figure and its reflection across the vertical axis: each
vertex is reflected square by square, and left and right are swapped.}
\end{omfigure}

\begin{proposition}[What a reflection preserves]\label{prop:g6:symmetry:preserved}
The symmetric of a figure across a line is a figure of the \emph{same
shape and same size}: lengths, angles, perimeters and areas are all
preserved. The symmetric of a segment is a segment, of a circle a
circle of the same radius, of a line a line.
\end{proposition}
\admitted

\begin{remark}
One thing is \emph{not} preserved: orientation. A reflected ``b''
becomes a ``d'' --- mirror writing. If tracing your figure and flipping
the tracing paper reproduces the image, the reflection is correct.
\end{remark}

\section{Axes of symmetry of a figure}

\begin{definition}[Axis of symmetry]\label{def:g6:symmetry:axis}
A line $d$ is an \emph{axis of symmetry}\index{axis of symmetry} of a
figure when the reflection across $d$ sends the figure exactly onto
itself --- folding along $d$ makes the two halves match.
\end{definition}

\begin{omfigure}
\begin{tikzpicture}[scale=0.62]
  % isosceles triangle
  \draw[thick] (0,0) -- (2.4,0) -- (1.2,2.6) -- cycle;
  \draw[omThm, dashed] (1.2,-0.5) -- (1.2,3.0);
  \node[below, font=\small] at (1.2,-0.7) {isosceles: $1$ axis};
  % rectangle
  \begin{scope}[xshift=5.6cm]
  \draw[thick] (0,0.3) rectangle (3.2,2.3);
  \draw[omThm, dashed] (1.6,-0.4) -- (1.6,3.0);
  \draw[omThm, dashed] (-0.5,1.3) -- (3.7,1.3);
  \node[below, font=\small] at (1.6,-0.7) {rectangle: $2$ axes};
  \end{scope}
  % square
  \begin{scope}[xshift=12.0cm]
  \draw[thick] (0.3,0.1) rectangle (2.7,2.5);
  \draw[omThm, dashed] (1.5,-0.5) -- (1.5,3.1);
  \draw[omThm, dashed] (-0.3,1.3) -- (3.3,1.3);
  \draw[omThm, dashed] (0.0,-0.2) -- (3.0,2.8);
  \draw[omThm, dashed] (0.0,2.8) -- (3.0,-0.2);
  \node[below, font=\small] at (1.5,-0.7) {square: $4$ axes};
  \end{scope}
\end{tikzpicture}

{\small Axes of symmetry of familiar shapes. Careful with the
rectangle: its diagonals are \emph{not} axes of symmetry (fold along a
diagonal --- the halves do not match unless it is a square).}
\end{omfigure}

\begin{example}\label{ex:g6:symmetry:shapes}
Counting axes of symmetry:
\begin{itemize}
  \item isosceles triangle: $1$ (through the top vertex and the middle
        of the base); equilateral triangle: $3$;
  \item rectangle: $2$ (the two ``middle lines'', \emph{not} the
        diagonals); rhombus: $2$ (its diagonals!); square: $4$;
  \item circle: every line through the center --- infinitely many.
\end{itemize}
\end{example}

\begin{example}[Symmetry proves equalities]\label{ex:g6:symmetry:proves}
Why are the two base angles of an isosceles triangle equal? Fold the
triangle along its axis of symmetry: the left half lands exactly on the
right half, so the left base angle lands on the right base angle ---
they must have the same measure
(\cref{prop:g6:symmetry:preserved}: reflections preserve angles).
\end{example}

\section{Perpendicular bisector}

\begin{definition}[Perpendicular bisector]\label{def:g6:symmetry:bisector}
The \emph{perpendicular bisector}\index{perpendicular bisector} of a
segment $[AB]$ is the line perpendicular to $[AB]$ through its
midpoint --- it is the axis of symmetry that swaps $A$ and $B$.
\end{definition}

\begin{proposition}[Equidistance]\label{prop:g6:symmetry:bisector}
A point $M$ is on the perpendicular bisector of $[AB]$ exactly when it
is at the same distance from $A$ and from $B$: $MA = MB$.
\end{proposition}
\admitted

\begin{method}[Compass construction]\label{met:g6:symmetry:bisector}
To draw the perpendicular bisector of $[AB]$ without measuring:
\begin{enumerate}
  \item open the compass to more than half of $AB$;
  \item draw two arcs centered at $A$, above and below the segment;
  \item with the \emph{same} opening, draw two arcs centered at $B$;
  \item join the two crossing points: this line is the perpendicular
        bisector (each crossing point is equidistant from $A$ and $B$).
\end{enumerate}
\end{method}

\section{Exercises}

\begin{exercise}[$\star$]\label{exo:g6:symmetry:1}
Copy on grid paper: a vertical axis $d$, a point $A$ two squares to its
left, a point $B$ five squares to its left and one square higher.
Construct their symmetric points $A'$ and $B'$ across $d$.
\end{exercise}

\begin{exercise}[$\star$]\label{exo:g6:symmetry:2}
Draw a line $d$ and a point $M$ at $3$ cm from $d$ (not on grid
paper). Construct the symmetric $M'$ of $M$ across $d$ with set square
and ruler, following \cref{met:g6:symmetry:construct}. How far apart
are $M$ and $M'$?
\end{exercise}

\begin{exercise}[$\star$]\label{exo:g6:symmetry:3}
Which capital letters of the alphabet, written in their simplest form,
have a vertical axis of symmetry (like A)? A horizontal one (like B)?
Both?
\end{exercise}

\begin{exercise}[$\star$]\label{exo:g6:symmetry:4}
How many axes of symmetry has: an equilateral triangle? a rectangle
(that is not a square)? a rhombus (not a square)? a square? a circle?
Draw each figure with its axes.
\end{exercise}

\begin{exercise}[$\star$]\label{exo:g6:symmetry:5}
True or false? ``The diagonals of a rectangle are axes of symmetry of
the rectangle.'' Explain with a folding argument or a drawing.
\end{exercise}

\begin{exercise}[$\star$]\label{exo:g6:symmetry:6}
A triangle $ABC$ is reflected across a line $d$ into $A'B'C'$. Given
$AB = 4$ cm, $\widehat{BAC} = 70^\circ$ and the perimeter of $ABC$ is
$12$ cm: give, without any construction, $A'B'$,
$\widehat{B'A'C'}$ and the perimeter of $A'B'C'$.
\end{exercise}

\begin{exercise}[$\star$]\label{exo:g6:symmetry:7}
Draw a segment $[AB]$ of $7$ cm and construct its perpendicular
bisector with the compass (\cref{met:g6:symmetry:bisector}). Choose any
point $M$ on it and check with the ruler that $MA = MB$.
\end{exercise}

\begin{exercise}[$\star\star$]\label{exo:g6:symmetry:8}
Two trees stand at points $A$ and $B$ of a (flat) garden. Where can one
plant a fountain so that it is at the same distance from both trees?
Describe \emph{all} the possible spots.
\end{exercise}

\begin{exercise}[$\star\star$]\label{exo:g6:symmetry:9}
On grid paper, draw the axis $d$ (a diagonal of the grid at
$45^\circ$) and a small L-shaped figure on one side. Construct its
symmetric figure across $d$. (Reflect each vertex: for a $45^\circ$
axis, a point $2$ squares right of the axis lands $2$ squares above
it.)
\end{exercise}

\begin{exercise}[$\star\star$]\label{exo:g6:symmetry:10}
Draw a circle of center $O$, a line $d$ through $O$, and a point $M$ on
the circle. Where is the symmetric of $M$ across $d$? Explain why the
reflection sends the circle onto itself.
\end{exercise}

\begin{exercise}[$\star\star$]\label{exo:g6:symmetry:11}
Using the equidistance property (\cref{prop:g6:symmetry:bisector}),
explain why the two crossing points of the arcs in
\cref{met:g6:symmetry:bisector} really do lie on the perpendicular
bisector of $[AB]$.
\end{exercise}

\begin{exercise}[$\star\star\star$]\label{exo:g6:symmetry:12}
A billiard ball at point $A$ must bounce off a straight wall $d$ and
reach point $B$ (both on the same side of the wall). Reflect $B$
across $d$ into $B'$, and draw the segment $[AB']$. Explain why the
best bouncing point is where $[AB']$ crosses the wall. (Idea: the path
$A \to P \to B$ has the same length as $A \to P \to B'$.)
\end{exercise}
